\scanpage{8}
\[
\operatorname{Hom}(Q,\operatorname{Hom}_A(M,P))\handwritten{\simeq}\operatorname{Hom}_A(M,\operatorname{Hom}(Q,P))\quad,
\]
(isomorphisme fonctoriel en l'objet variable $Q$ de $C_1$), où ce fois-ci $M$ \uline{et}
$P$ sont des $A$-modules à gauche ($P$ à valeurs dans $C_1$). On utilisera sans
référence spéciale les sorites duaux de ceux donnés pour $\otimes$, qui sont une
simple traduction de ces derniers. L'intérêt du traitement abst\deleted{\unsure{a}}rait donné
ici est précisément, entre autres, d'unifier certains énoncés \inserted{\deleted{\unsure{des}}\handwritten{sur \unsure{les|des}} modules habituels} portant, \deleted{\unsure{l'u}}
l'un sur des produits tensoriels, l'autre sur des $\operatorname{Hom}$,
(il faut avoir vu bien entendu qu'on retroyve \deleted{c}\handwritten{l}es \deleted{derniers lorsque}\inserted{\handwritten{\unsure{objets} \unsure{traditionnels|habituel} si}}
$C'$ resp. $C_1$ est la catégorie $(\mathrm{Ab})$, i.e. comme cas particulier de produits
tensoriels avec des $A$-$(\mathrm{Ab})$-modules et des $A$-$(\mathrm{Ab})^{o}$-modules\handwritten{)}.

Ces rappels étant faits, par définition la donnée d'un homomorphisme
$(\handwritten{*})$ équivaut à la donnée d'un $A$-homomorphisme
\[
M\handwritten{\longrightarrow}\operatorname{Hom}(T(A),T(M))\quad,
\]
or un tel s'obtient de façon naturelle en composant
\[
M\handwritten{\longrightarrow}\operatorname{Hom}_A(A,M)\handwritten{\longrightarrow}\operatorname{Hom}(T(A),T(M))
\]
Ceci posé, on a:

\uline{Proposition} 2.6. L'homomorphisme $(\handwritten{*}):T(A)\otimes_A M\handwritten{\to}T(M)$ est fonctoriel.
Pour que ce soit un isomorphisme, il f et s que $T$ soit\handwritten{\unsure{/}}exact à droite.

\deleted{Suffisance}\inserted{Nécessité} triviale, pour la suffisance, on note que tout homomorphisme \deleted{de} $T'(M)\handwritten{\to}T(M)$ de foncteurs exacts à droite, qui induit un isomorphisme $T'(A)\handwritten{\to}T(A)$, est un isomorphisme \deleted{pour} (car c'est un isomorphisme
pour un argument $L$ libre de type fini, donc aussi pour un conoyau
$L_1\handwritten{\to}L_{\unsure{o|0}}$ de tels arguments, puisque $T,T'$ sont exacts à droite ...).

\uline{Corollaire} 2.7. Le foncteur $T\handwritten{\leadsto}T(A)$ est une équivalence de la catégorie
\deleted{des} $\underline{\operatorname{Hom}}(\underline C,\underline C')_{\deleted{g}d}$ des foncteurs exacts à \deleted{gauche}\inserted{droite} de $\underline C$ \handwritten{d}ans $\underline C'$, dans\handwritten{\unsure{/}}la
catégorie $A(\deleted{\unsure{\mathrm{Ab}}}\handwritten{\underline C'})$ des $A$-\inserted{\handwritten{$\underline C'$-}}modules quelconques, une équivalence\handwritten{\unsure{/}}quasi-inverse
est le foncteur $H\handwritten{\leadsto}T_H$, où $T_H(M)=H\otimes_A M$.

\uline{Remarque} 2.8. L'homomorphisme $(\handwritten{*})$ est défini également, sans hypothèse
noethé\unsure{ri|fi}enne sur $A$, lorsque $\underline C$ est la catégorie de tous les $A$-modules,
p\handwritten{\unsure{our}}vu que dans $\underline C'$ les sommes directes quelconques existent (donc que
le premier membre ait un sens). Alors 2.6. \deleted{peut} reste valable en y
précisant: $T$ exact à droite \uline{et} commute aux sommes directes ; et de même
pour 2.7.

Nous auronségalement à utiliser la situation duale d'un foncteur
% archive sheet number: 7

